\documentclass[11pt]{article}
\usepackage[margin=1in]{geometry}
\usepackage{amsmath,amssymb,amsthm,mathtools}
\usepackage[T1]{fontenc}
\usepackage[expansion=false]{microtype}
\usepackage[hidelinks]{hyperref}
\newtheorem{theorem}{Theorem}[section]
\newtheorem{lemma}[theorem]{Lemma}
\newtheorem{proposition}[theorem]{Proposition}
\newcommand{\ddc}{\mathrm d\mathrm d^c}
\newcommand{\PSH}{\operatorname{PSH}}
\newcommand{\CP}{\mathbb P}
\newcommand{\vol}{\mathrm dV}
\newcommand{\norm}[1]{\left\|#1\right\|}
\title{Curvature pairs and the full determinant endpoint on elliptic threefolds}
\author{}
\date{}
\begin{document}
\maketitle
\begin{abstract}
A sign condition on the two real curvature forms of a principal elliptic
bundle over a K\"ahler surface gives an explicit determinant estimate over
all normalized positive Gauduchon currents, including singular and degenerate
currents. A local bounded-potential energy identity proves the exact lower
Monge--Amp\`ere mass of a split metric. A matrix calculation transports the
endpoint to every Hermitian background. Under the stated sign condition,
\(Q\not\equiv0\) excludes every pluriclosed metric, so the criterion extends beyond the mechanism requiring
the first two metric powers to be \(\ddc\)-closed. Concrete examples include
negative-line-bundle quotients over \(\CP^1\times\CP^1\) and a complex structure
on \(S^3\times S^3\) admitting no LCK metric. Primitive curvature on the base
has the opposite sign. The result is restricted to the stated threefold
geometry; higher dimensions and general elliptic lower-mass collapse are
not resolved here.
\end{abstract}

\section{Framework and the geometric statement}

Use \(\mathrm d^c=(i/(4\pi))(\bar\partial-\partial)\) and
\(\ddc=(i/(2\pi))\partial\bar\partial\), with exterior powers without
factorials. All quasi-psh potentials are genuine canonical upper
representatives. For a Hermitian form \(h\) on a compact complex threefold,
put
\[
 v_-(h)=\inf_{u\in\PSH(X,h)\cap L^\infty(X)}
                     \int_X(h+\ddc u)^3.
\]
\[
 {\cal G}_h=\{T\ge0:\deg T=(2,2),\
             \ddc T=0,\ \int_Xh\wedge T=1\}.
\]
The bounded products are the local Bedford--Taylor products with smooth
nonclosed-background corrections \cite[Section 2.1]{BGL}.

Fix a smooth positive probability \(\vol\). Write \(\sigma_h=h\wedge T\).
In an \(h\)-unitary frame the positive matrix density \(P_h\), with
\(\operatorname{tr}P_h=3\), is specified by
\[
 \beta\wedge T=\tfrac13\operatorname{tr}(BP_h)\,\sigma_h,
 \qquad
 \Gamma_0(T)=\left(\frac{\vol}{h^3}\right)^{1/3}
                   \det(P_h)^{1/3}\sigma_h,
 \qquad k_T=\frac{\mathrm d\Gamma_0(T)}{\vol}.
\]
Here \(B\) is the relative matrix of a smooth real \((1,1)\)-form.
We use as an explicit input the critical determinant/lower-mass duality
of the underlying framework:
\begin{equation}\label{eq:duality}
 v_-(h)>0\ \Longleftrightarrow\
 \sup_{T\in{\cal G}_h}\norm{k_T}_{L^{3/2}(\vol)}<\infty,
 \qquad
 \sup_{T\in{\cal G}_h}\norm{k_T}_{L^{3/2}(\vol)}
                       =v_-(h)^{-1/3}.
\end{equation}
Its current family includes singular and degenerate tests, and
\(\Gamma_0(T)\) is absolutely continuous as specified in that framework.
The geometric and bounded-potential conclusions below establish its
antecedent; no estimate on one maximizing test is substituted for it.

Let \(B\) be a compact connected K\"ahler surface with K\"ahler form
\(\omega\). Let \(\pi:X\to B\) be a holomorphic principal elliptic bundle
with fiber \(E=\mathbb C/\Lambda\). Suppose there is an invariant smooth
\((1,0)\)-form \(\eta\), restricting to \(\mathrm dz\) on every fiber,
such that
\[
 \mathrm d\eta=\pi^*F,\qquad F=\alpha+i\beta,
\]
where \(\alpha,\beta\) are smooth closed real \((1,1)\)-forms on \(B\).
Set
\[
 Q=\alpha^2+\beta^2,\qquad
 e=i\eta\wedge\bar\eta,\qquad g=\pi^*\omega+s e,\quad s>0,
 \qquad A_E=\int_E i\,\mathrm dz\wedge\mathrm d\bar z.
\]

\begin{theorem}\label{thm:bundle}
Assume \(Q\ge0\) as a top-degree form on \(B\), and \(Q\not\equiv0\).
Then \(g\) is Hermitian and
\[
 V:=\int_Xg^3=3s A_E\int_B\omega^2,\qquad v_-(g)=V.
\]
Every bounded genuine \(g\)-psh \(u\) has the exact mass formula
\begin{equation}\label{eq:bundle-mass}
 \int_X(g+\ddc u)^3
 =V+\frac{3s}{2\pi}\int_X E_u\wedge\pi^*Q\ge V,
 \qquad E_u=\tfrac12\ddc(u^2)-u\ddc u\ge0.
\end{equation}
For every Hermitian \(h\ge c g\), \(c>0\),
\begin{equation}\label{eq:endpoint}
 \sup_{T\in{\cal G}_h}\norm{k_T}_{L^{3/2}(\vol)}
       \le c^{-1}V^{-1/3},\qquad v_-(h)\ge c^3 V.
\end{equation}
In particular these bounds apply to every Hermitian metric on \(X\),
with its own positive comparison constant. No pluriclosed Hermitian
metric exists on \(X\).
\end{theorem}

The signed-torsion lower-mass lemma used in the proof is already present
in Angella--Guedj--Lu \cite[Theorem 3.2]{AGL}. Their Vaisman threefold
result is also known \cite[Proposition 3.10]{AGL}. We give the bounded
local proof, the full-current change of background, and explicit
curvature-pair applications. No literature-priority assertion is made.
The hypotheses do not require bounded upper mass, a bound on potential
oscillations, or a pluriclosed metric.

\section{Bounded energy and the cubic lower mass}

\begin{lemma}\label{lem:energy}
For every bounded genuine quasi-psh \(u\), the signed closed products
\(C_r=(\ddc u)^r\) satisfy
\(\ddc(u C_r)=C_{r+1}\), for \(r=0,1,2\).
The current \(E_u=\frac12\ddc(u^2)-u\ddc u\) is positive.
\end{lemma}
\begin{proof}
Locally write \(u=p-a\), with bounded psh \(p\) and smooth real \(a\).
Define \(C_r\) by expanding \(\ddc p-\ddc a\).
The recursive products, closedness, multilinearity and decreasing
continuity are precisely \cite[(2.1), Theorem 2.1, Corollary 2.2]{BT}.
The weighted continuity needed here is
\cite[Theorem 2.4]{BT}. Multiplying by \(p-a\) and applying \(\ddc\)
inserts \(\ddc p-\ddc a\); expansion and Pascal's identity prove the
recursion. The local definitions agree under a change of smooth
decomposition, hence define global order-zero currents.

On a smaller ball let smooth bounded psh \(p_j\downarrow p\), by local
convolution, and put \(u_j=p_j-a\). Local uniform boundedness gives
\(u_j^2\to u^2\) in \(L^1\). Expand
\[
 u_j\ddc u_j
 =p_j\ddc p_j-a\ddc p_j-p_j\ddc a+a\ddc a.
\]
Weighted Bedford--Taylor continuity handles the first term;
distributional convergence and dominated convergence handle the others.
Consequently
\[
 \mathrm du_j\wedge\mathrm d^c u_j
 =\tfrac12\ddc(u_j^2)-u_j\ddc u_j\longrightarrow E_u.
\]
Each smooth term is positive. Thus the weak limit is a positive
order-zero current, and the intrinsic formula glues it globally.
No global smooth approximation preserving the Hermitian background
has been used.
\end{proof}

\begin{proposition}\label{prop:cubic}
If a Hermitian form \(g\) on a compact complex threefold satisfies
\(\ddc(g^2)=0\) and \(S=-\ddc g\ge0\), then
\[
 \int_X(g+\ddc u)^3=\int_Xg^3+3\int_XE_u\wedge S
\]
for every bounded genuine \(g\)-psh \(u\). In particular
\(v_-(g)=\int_Xg^3\).
\end{proposition}
\begin{proof}
The bounded-product expansion is
\[
 (g+\ddc u)^3=g^3+3g^2 C_1+3g C_2+C_3.
\]
Lemma~\ref{lem:energy} and distributional Stokes give
\[
 \int_X C_3=0,\qquad
 \int_X g^2 C_1=\int_Xu\ddc(g^2)=0,\qquad
 \int_X g C_2=\int_Xu C_1\ddc g.
\]
Since \(\ddc g\) is closed, the definition of \(E_u\) gives
\[
 \int_X u C_1\ddc g=-\int_XE_u\ddc g=\int_XE_u S.
\]
This proves the identity and positivity. The constant \(u=0\) attains
the lower bound. The nonclosed product convention and positivity for
bounded potentials are those of \cite[Section 2.1, Proposition 2.1]{BGL};
the BMP hypothesis in their following subsection is not needed here.
\end{proof}

If \(S\not\equiv0\) and a positive Hermitian \(b\) were pluriclosed, then
\[
 0=\int_Xg\wedge\ddc b=\int_X b\wedge\ddc g
                   =-\int_Xb\wedge S<0.
\]
A nonzero smooth positive \((2,2)\)-form has strictly positive pairing
with \(b\) on an open set. This excludes every pluriclosed metric, hence
every metric satisfying both identities \(\ddc b=\ddc(b^2)=0\).

\section{Change of background for all currents}

The matrix density definition is independent of the auxiliary Hermitian
form. Work in an \(h\)-unitary frame, with \(G\) the relative positive
matrix of \(g\), and \(P=P_h\). Put
\[
 a=\tfrac13\operatorname{tr}(GP),\qquad
 \sigma_g=a\sigma_h,\qquad P_g=a^{-1}G^{1/2}PG^{1/2}.
\]
Then \(\operatorname{tr}P_g=3\). A form with \(h\)-matrix \(B\) has
\(g\)-matrix \(G^{-1/2}BG^{-1/2}\), which verifies the trace pairing.
Moreover
\[
 \det P_g=a^{-3}\det G\det P,\qquad g^3=\det G\,h^3.
\]
Substitution into \(\Gamma_0(T)\) cancels both \(a\) and \(\det G\).
The calculation holds \(\sigma_h\)-almost everywhere for the
Radon matrix density of any positive current. Since \(G>0\) and
\(\operatorname{tr}P=3\), \(a>0\); no invertibility of \(P\), smooth
coefficient density or removal of a singular trace component is needed.
For a constant \(m>0\) one also has \(k_{mT}=m k_T\).

Take any \(T\in{\cal G}_h\). If \(h\ge c g\), then
\[
 0<m:=\int_Xg\wedge T\le c^{-1},\qquad U=T/m\in{\cal G}_g.
\]
Proposition~\ref{prop:cubic} and the explicit input
\eqref{eq:duality} give
\[
 \norm{k_T}_{3/2}=m\norm{k_U}_{3/2}
                         \le c^{-1}\left(\int_Xg^3\right)^{-1/3}.
\]
The same duality gives the all-bounded-potential lower mass for \(h\).
This argument treats the full normalized-current family.

\section{Curvature computation and primitive exclusion}

Suppress pullbacks. The type and closedness assumptions give
\(\partial\eta=0\), \(\bar\partial\eta=F\),
\(\partial\bar\eta=\bar F\) and \(\bar\partial\bar\eta=0\).
Thus
\[
 \partial e=-i\eta\wedge\bar F,\qquad
 \bar\partial e=iF\wedge\bar\eta,\qquad
 \ddc e=-\frac1{2\pi}F\wedge\bar F=-\frac1{2\pi}Q.
\]
Since \(e^2=0\) and the base has complex dimension two,
\[
 \ddc g=-\frac{s}{2\pi}\pi^*Q,\qquad
 \ddc(g^2)=-\frac{s}{\pi}\pi^*(\omega\wedge Q)=0,\qquad
 g^3=3s\,\pi^*(\omega^2)\wedge e.
\]
Positivity on the horizontal quotient and on the vertical kernel proves
that \(g\) is Hermitian. Fiber integration gives \(V\) and completes
Theorem~\ref{thm:bundle}.

Balancedness of this split metric is equivalent to \(\omega\wedge F=0\):
\(\mathrm d(g^2)=2s\,\omega\wedge\mathrm de\), and its two independent
vertical types force the equivalence. Hence a primitive pair satisfies
\(\omega\wedge\alpha=\omega\wedge\beta=0\). Diagonalizing a primitive
real \((1,1)\)-form on a surface gives eigenvalues \(\lambda,-\lambda\),
so
\[
 \alpha^2=-\tfrac12|\alpha|_\omega^2\omega^2,\qquad
 Q=-\tfrac12(|\alpha|_\omega^2+|\beta|_\omega^2)\omega^2.
\]
For a nonzero primitive pair the sign required by
Theorem~\ref{thm:bundle} fails somewhere. In the primitive line-bundle
model \(F=2\pi i c\), this is \(Q=4\pi^2c^2\le0\).
The mass identity then has the opposite energy sign. The sufficient
condition excludes that degeneration; it does not determine the
unrestricted primitive-bundle infimum.

\section{Negative-line-bundle examples beyond scalar Hopf}

Let \(B=\CP^1\times\CP^1\) with normalized Fubini--Study forms
\(\omega_1,\omega_2\), and set
\[
 \Omega=a\omega_1+b\omega_2,\quad a,b\in\mathbb N_{>0},\qquad
 L={\cal O}(-a,-b),\qquad
 X=(L^\times)/\langle v\mapsto\rho v\rangle,\quad0<\rho<1.
\]
Use the tensor-product tautological norm, with squared local frame norm
\((1+|z_1|^2)^a(1+|z_2|^2)^b\). For
\(r=\log\|v\|^2\),
\[
 \ddc r=\pi^*\Omega,\qquad\eta=\partial r,\qquad
 \mathrm d\eta=2\pi i\,\pi^*\Omega.
\]
The form descends because multiplication changes \(r\) by a constant.
The quotient is a holomorphic elliptic bundle with fiber
\(\mathbb C^*/\rho^{\mathbb Z}\). Freeness, properness and compactness
follow from the nonzero radial period and the compact unit-circle bundle
with a fundamental radial interval. Thus
for the split metric \(g=\pi^*\Omega+s i\eta\wedge\bar\eta\),
\[
 Q=4\pi^2\Omega^2>0,\quad
 A_E=4\pi\log(1/\rho),\quad \int_B\Omega^2=2ab,\quad
 v_-(g)=24\pi s ab\log(1/\rho).
\]
The corresponding bound \eqref{eq:endpoint} holds for every \(h\).

These examples are not scalar Hopf threefolds. Put
\(A=\eta+\bar\eta\), \(B_1=i(\eta-\bar\eta)\).
Then \(\mathrm dA=0\) and
\(\mathrm dB_1=-4\pi\pi^*\Omega\).
If \(\pi^*(\omega_1-\omega_2)=\mathrm d\zeta\), average \(\zeta\)
over the compact elliptic group. An invariant one-form has the unique
decomposition
\(\zeta=\pi^*\xi+(\pi^*f)A+(\pi^*j)B_1\).
Evaluation on one vertical and one horizontal vector gives
\(\mathrm df=\mathrm dj=0\). The resulting basic equation is
\[
 \omega_1-\omega_2=\mathrm d\xi-4\pi j\Omega.
\]
Integration over the two factor spheres forces
\(1=-4\pi ja\) and \(-1=-4\pi jb\), a contradiction.
Thus \(H^2(X,\mathbb R)\ne0\), whereas a scalar Hopf threefold is
diffeomorphic to \(S^5\times S^1\) and has \(H^2=0\).
The exact nonzero semipositive form \(\pi^*\Omega\) also excludes every
balanced metric: its pairing with a closed positive metric square would
be both zero by Stokes and strictly positive.

\section{A supplementary Calabi--Eckmann example with no LCK metric}

This construction is a supplementary instance of
Theorem~\ref{thm:bundle}. Let \(\tau=A+iB\), \(A<0\), \(B>0\), and
take the holomorphic \(\mathbb C\)-action
\[
 t\cdot(z,w)=(e^t z,e^{\tau t}w)
 \quad\hbox{on }(\mathbb C^2\setminus\{0\})^2.
\]
Put \(r_1=\log|z|^2\), \(r_2=\log|w|^2\).
Each orbit meets \(|z|=|w|=1\) at exactly one point, using
\[
 \Re t=-r_1/2,\qquad \Im t=(r_2-A r_1)/(2B).
\]
These formulas give a smooth product trivialization by
\(\mathbb C\times(S^3\times S^3)\); hence the action is free and proper.
Local holomorphic slices set a nonzero \(z\)-coordinate to one using a
branch of its logarithm, with coordinates
\(z_{\rm other}/z_{\rm chosen}\) and \(w_j z_{\rm chosen}^{-\tau}\).
Their transitions are holomorphic. The quotient \(X_\tau\) is therefore
a compact complex threefold diffeomorphic to \(S^3\times S^3\).
The slice coordinates use local logarithm branches; their residual
periods are accounted for in the elliptic fibers below.

Projectivization gives an elliptic bundle over \(\CP^1\times\CP^1\).
The form
\[
 \eta=\tau\partial r_1-\partial r_2
\]
is invariant and annihilates the holomorphic action generator, so it
descends. On a fiber it is \(\mathrm d\chi\), with
\(\chi=\tau\log v_1-\log v_2\) and lattice
\(\langle2\pi i\tau,-2\pi i\rangle\).
Consequently \(A_E=8\pi^2B\). With \(\omega=\omega_1+\omega_2\),
\[
 F=2\pi i(\tau\omega_1-\omega_2),\quad
 \alpha=-2\pi B\omega_1,\quad
 \beta=2\pi(A\omega_1-\omega_2),\quad
 Q=-8\pi^2A\,\omega_1\wedge\omega_2>0.
\]
The split metric therefore has
\[
 v_-(g)=48\pi^2sB,\qquad
 \sup_{T\in{\cal G}_h}\norm{k_T}_{3/2}
                         \le c^{-1}(48\pi^2sB)^{-1/3}
 \quad(h\ge c g).
\]

For every parameter in this range, the equation
\[
 \mathrm d\Re\eta=-2\pi B\,\pi^*\omega_1
\]
makes the nonzero semipositive form \(\pi^*\omega_1\) exact. It excludes
every balanced metric and hence every K\"ahler metric. The quotient is simply connected,
so every closed real one-form is exact. Any LCK metric would therefore
be globally conformal to a K\"ahler metric, a contradiction.
Thus every example in this range admits no LCK metric, and in particular no Vaisman
metric. The nonzero \(Q\) also excludes every pluriclosed metric.
This family supplies nonclosed endpoint examples beyond both the
old identities and the scalar Hopf or Vaisman geometries.

\section{Scope}

The conclusions are quantitative for all normalized currents and all
bounded admissible potentials on the stated elliptic threefolds.
They do not classify all positive-lower manifolds or solve higher
dimensions, changing-sign curvature, or primitive elliptic bundles,
including the curve-free K3 problem. The signed-torsion lemma and the
input duality are explicitly attributed or stipulated; the examples
and change-of-background calculation give the claimed geometric
sufficient conditions. No estimate on a selected adjoint family or
publication-priority claim replaces those conclusions.

\begin{thebibliography}{9}
\bibitem{BT}
E. Bedford and B. A. Taylor, \emph{A new capacity for plurisubharmonic
functions}, Acta Math. \textbf{149} (1982), 1--40, Section 2.
\href{https://www.math.stonybrook.edu/~ebedford/PapersForM655/BTNewCapacity.pdf}
{Author's university copy}.
\bibitem{BGL}
S. Boucksom, V. Guedj and C. H. Lu, \emph{Volumes of Bott--Chern classes},
arXiv:2406.01090v1 (3 June 2024), Section 2.1.
\href{https://arxiv.org/abs/2406.01090v1}{Official primary manuscript}.
\bibitem{AGL}
D. Angella, V. Guedj and C. H. Lu, \emph{Plurisigned hermitian metrics},
arXiv:2207.04705v1 (11 July 2022), Theorem 3.2 and Proposition 3.10;
published in Trans. Amer. Math. Soc. \textbf{376} (2023), 4631--4659.
\href{https://arxiv.org/abs/2207.04705v1}{Primary preprint}.
\end{thebibliography}
\end{document}
